\section{Anexo - Atributos de calidad}

La arquitectura y el diseño de \textbf{SocioUnido} fueron concebidos bajo lineamientos rigurosos de ingeniería de software, priorizando un conjunto de atributos de calidad que responden de manera directa a las necesidades críticas del ecosistema deportivo. A continuación, se detallan los pilares fundamentales que guían el desarrollo de la plataforma:

\subsection{Seguridad}

Representa uno de los pilares centrales del sistema. Dado que la plataforma gestiona información altamente sensible (datos personales, registros financieros y credenciales de acceso), se aplican estrictas medidas de protección para garantizar la confidencialidad, la privacidad y la integridad absoluta de los datos, resguardando tanto a los clubes como a su masa societaria frente a posibles vulnerabilidades.

\subsection{Escalabilidad}

La arquitectura de la plataforma tiene el propósito fundamental de priorizar la adaptabilidad ante variaciones de carga. Al adoptar un diseño modular basado en microservicios, el sistema permite escalar cada componente de forma independiente según la demanda específica de recursos. Esta separación de responsabilidades facilita alcanzar cotas de rendimiento sumamente altas sin comprometer la estabilidad global de la aplicación.

\subsection{Disponibilidad y resiliencia}

Debido al flujo continuo de personas en las instalaciones del club y a la existencia de momentos operativos críticos (como los días de partido o eventos de asistencia masiva), el sistema exige un tiempo de actividad (\textit{uptime}) máximo. Se diseñaron estrategias de mitigación para cortar posibles vectores de ataque y aislar problemas técnicos, asegurando que un fallo aislado no derive en una interrupción total del servicio.

\subsection{Rendimiento y eficiencia}

Si bien la seguridad y la disponibilidad son los núcleos de la arquitectura, el rendimiento es un factor clave para garantizar una adopción exitosa. Se optimizó el consumo de recursos y los tiempos de respuesta para ofrecer una interacción fluida y ágil, logrando que el uso diario de la aplicación resulte una experiencia agradable y no represente un obstáculo burocrático o tecnológico para el usuario.

\subsection{Usabilidad y accesibilidad}

En estrecha relación con el rendimiento, la interfaz y los flujos de trabajo fueron diseñados para reducir al mínimo la fricción operativa. Para alcanzar este estándar, se llevaron a cabo exhaustivas reuniones de validación de producto, se realizaron encuestas y se procesó un volumen significativo de \textit{feedback} provisto por los propios usuarios. Esto permitió pulir el sistema para que sea intuitivo y accesible frente a una demografía con distintos niveles de alfabetización digital.

\newpage

\subsection{Interoperabilidad}

La capacidad del sistema para integrarse con actores y sistemas externos se centró estratégicamente en dos frentes de expansión futura. En primer lugar, la aplicación está estructurada para facilitar la interacción directa con modelos de inteligencia artificial y agentes autónomos; ante el actual superciclo de innovación tecnológica, resulta indispensable contar con interfaces preparadas por si esta modalidad se consolida como el estándar masivo de interacción. En segundo lugar, el diseño contempla la futura integración a nivel de \textit{hardware} (como la conexión con molinetes y controles de acceso físicos) para agilizar y blindar de forma definitiva el ingreso a los recintos deportivos.
